% ---------------------------------------------------------------------------
% ---------------------------------------------------------------------------
% Junção de templates encontrados no Overleaf
% facilitando a vida do aluno de pós da UFABC
% Modelo LaTex para preparação do documento final de Dissertação de Mestrado
% ninguém do programa de pós da informação validou se presta
% ---------------------------------------------------------------------------
% ---------------------------------------------------------------------------

\documentclass[
	% -- opções da classe memoir --
	12pt,					% tamanho da fonte
	openright,				% capítulos começam em pág ímpar (insere página vazia caso preciso)
	twoside,					% para impressão em verso e anverso. Oposto a oneside
	a4paper,					% tamanho do papel. 
	% -- opções da classe abntex2 --
	%chapter=TITLE,			% títulos de capítulos convertidos em letras maiúsculas
	%section=TITLE,			% títulos de seções convertidos em letras maiúsculas
	%subsection=TITLE,		% títulos de subseções convertidos em letras maiúsculas
	%subsubsection=TITLE,	% títulos de subsubseções convertidos em letras maiúsculas
	% -- opções do pacote babel --
	english,					% idioma adicional para hifenização
	%french,					% idioma adicional para hifenização
	%spanish,				% idioma adicional para hifenização
	brazil					% o último idioma é o principal do documento
	]{abntex2}

% ---------------------
% Pacotes OBRIGATÓRIOS
% ---------------------
\usepackage{lmodern}				% Usa a fonte Latin Modern			
\usepackage[T1]{fontenc}			% Selecao de codigos de fonte.
\usepackage[utf8]{inputenc}		% Codificacao do documento (conversão automática dos acentos)
\usepackage{lastpage}			% Usado pela Ficha catalográfica
\usepackage{indentfirst}			% Indenta o primeiro parágrafo de cada seção.
\usepackage{color}				% Controle das cores
\usepackage{graphicx,graphicx}	% Inclusão de gráficos
\usepackage{epsfig,subfig}		% Inclusão de figuras
\usepackage{microtype} 			% Melhorias de justificação
% ---------------------
		
% ---------------------
% Pacotes ADICIONAIS
% ---------------------
\usepackage{lipsum}						% Geração de dummy text
\usepackage{amsmath,amssymb,mathrsfs}	% Comandos matemáticos avançados 
\usepackage{setspace}  					% Para permitir espaçamento simples, 1 1/2 e duplo
\usepackage{verbatim}					% Para poder usar o ambiente "comment"
\usepackage{tabularx} 					% Para poder ter tabelas com colunas de largura auto-ajustável
\usepackage{afterpage} 					% Para executar um comando depois do fim da página corrente
\usepackage{url} 						% Para formatar URLs (endereços da Web)
% ---------------------

% ---------------------
% Pacotes de CITAÇÕES
% ---------------------
\usepackage[brazilian,hyperpageref]{backref}	% Paginas com as citações na bibl
\usepackage[alf]{abntex2cite}				% Citações padrão ABNT (alfa)
%\usepackage[num]{abntex2cite}				% Citações padrão ABNT (numericas)
% ---------------------

% Configurações de CITAÇÕES para abntex2
% --- 
% CONFIGURAÇÕES DE PACOTES
% --- 

% ---
% Configurações do pacote backref
% Usado sem a opção hyperpageref de backref
\renewcommand{\backrefpagesname}{Citado na(s) página(s):~}
% Texto padrão antes do número das páginas
\renewcommand{\backref}{}
% Define os textos da citação
\renewcommand*{\backrefalt}[4]{
	\ifcase #1 %
		Nenhuma citação no texto.%
	\or
		Citado na página #2.%
	\else
		Citado #1 vezes nas páginas #2.%
	\fi}%
% ---

% Inclusão de dados para CAPA e FOLHA DE ROSTO (título, autor, orientador, etc.)
% ---
% Informações de dados para CAPA e FOLHA DE ROSTO
% ---
\titulo{What Weee Are: jogo eletrônico com abordagem à reciclagem de resíduos eletrônicos}
\autor{Richard da Cruz Lopes}
\local{Santo André - SP}
\data{16 de Dezembro de 2022}
\orientador{Prof Dr André Luiz Brandão}
\instituicao{%
  Universidade Federal do ABC -- UFABC
  \par
  Centro de Matemática, Cognição e Computação
  \par
  Bacharelado em Ciência da Computação}
\tipotrabalho{Dissertação (Projeto de Graduação de Curso)}
% O preambulo deve conter o tipo do trabalho, o objetivo,
% o nome da instituição e a área de concentração
\preambulo{\textbf{Dissertação de Graduação} como parte dos requisitos necessários para a obtenção do Título de Bacharel em Ciência da Computação.}
% ---

% Inclui Configurações de aparência do PDF Final
%  Configurações de aparência do PDF final
% NÃO ALTERAR!!!

% alterando o aspecto da cor azul
\definecolor{blue}{RGB}{41,5,195}

% informações do PDF
\makeatletter
\hypersetup{
     	%pagebackref=true,
		pdftitle={\@title}, 
		pdfauthor={\@author},
    		pdfsubject={\imprimirpreambulo},
	    pdfcreator={LaTeX with abnTeX2},
		pdfkeywords={abnt}{latex}{abntex}{abntex2}{trabalho acadêmico}, 
		colorlinks=true,       		% false: boxed links; true: colored links
    		linkcolor=blue,          	% color of internal links
    		citecolor=blue,        		% color of links to bibliography
    		filecolor=magenta,      		% color of file links
		urlcolor=blue,
		bookmarksdepth=4
} 
\makeatother
% --- 

% O tamanho da identação do parágrafo é dado por:
\setlength{\parindent}{1.3cm}

% Controle do espaçamento entre um parágrafo e outro:
\setlength{\parskip}{0.2cm}  % tente também \onelineskip

% ---------------------
% Compila o indice
% ---------------------
\makeindex
% ---------------------

%%%%%%%%%%%%%%%%%%%%%%%%%%%
%%  INICIO DO DOCUMENTO  %%
%%%%%%%%%%%%%%%%%%%%%%%%%%%
\begin{document}

% Retira espaço extra obsoleto entre as frases.
\frenchspacing

% ----------------------------------------------------------
% ELEMENTOS PRÉ-TEXTUAIS (Capa, Resumo, Abstract, etc.)
% ----------------------------------------------------------
\pretextual

% Capa
% ---
% Impressão da Capa
% ---
  \begin{capa}%
    \begin{figure}[h!]%
        \centering%
        \includegraphics[scale=1.2]{figs/logo.png}%
      \end{figure}%
    \center
	\ABNTEXchapterfont\large{Universidade Federal do ABC \\ Centro de Matemática, Cognição e Computação \\ Bacharel em Ciência da Computação}
	%\vspace{1.5cm}

    \vfill
    \ABNTEXchapterfont\bfseries\LARGE\imprimirtitulo
    \vfill

	%\vfill
	\ABNTEXchapterfont\large\imprimirautor
	\vfill
%
	Número de Ordem : XXXX
	
    \large\imprimirlocal, \large\imprimirdata

    \vspace*{1cm}
  \end{capa}
% ---

% Folha de rosto (o * indica que haverá a ficha bibliográfica)
\imprimirfolhaderosto*

% Imprimir Ficha Catalografica
% ---
% Ficha Catalográfica
% ---
% Isto é um exemplo de Ficha Catalográfica, ou ``Dados internacionais de
% catalogação-na-publicação''. Você pode utilizar este modelo como referência. 
% Porém, talvez a biblioteca lhe fornece um PDF
% com a ficha catalográfica definitiva após a defesa do trabalho. Quando estiver
% com o documento, salve-o como PDF no diretório do seu projeto e substitua todo
% o conteúdo de implementação deste arquivo pelo comando abaixo:
%
% \begin{fichacatalografica}
%     \includepdf{fig_ficha_catalografica.pdf}
% \end{fichacatalografica}
\begin{fichacatalografica}
	\vspace*{\fill}					% Posição vertical
	\hrule							% Linha horizontal
	\begin{center}					% Minipage Centralizado
	\begin{minipage}[c]{12.5cm}		% Largura
	
	\imprimirautor
	
	\hspace{0.5cm} \imprimirtitulo  / \imprimirautor. --
	\imprimirlocal, \imprimirdata-
	
	\hspace{0.5cm} \pageref{LastPage} p. : il. (algumas color.) ; 30 cm.\\
	
	\hspace{0.5cm} \imprimirorientadorRotulo~\imprimirorientador\\
	
	\hspace{0.5cm}
	\parbox[t]{\textwidth}{\imprimirtipotrabalho~--~\imprimirinstituicao,
	\imprimirdata.}\\
	
	\hspace{0.5cm}
		1. Palavra-chave1.
		2. Palavra-chave2.
		I. Orientador.
		II. Universidade xxx.
		III. Faculdade de xxx.
		IV. Título\\ 			
	
	\hspace{8.75cm} CDU 02:141:005.7\\
	
	\end{minipage}
	\end{center}
	\hrule
\end{fichacatalografica}
% ---

% Inserir Folha de Aprovação
% ---
% Assinaturas
% ---
% Isto é um exemplo de Folha de aprovação, elemento obrigatório da NBR
% 14724/2011 (seção 4.2.1.3). Você pode utilizar este modelo até a aprovação
% do trabalho. Após isso, substitua todo o conteúdo deste arquivo por uma
% imagem da página assinada pela banca com o comando abaixo:
%
% \includepdf{folhadeaprovacao_final.pdf}
%
\begin{folhadeaprovacao}

  \begin{center}
    {\ABNTEXchapterfont\large\imprimirautor}

    \vspace*{\fill}\vspace*{\fill}
    \begin{center}
      \ABNTEXchapterfont\bfseries\Large\imprimirtitulo
    \end{center}
    \vspace*{\fill}
    
    \hspace{.45\textwidth}
    \begin{minipage}{.5\textwidth}
        \imprimirpreambulo
    \end{minipage}%
    \vspace*{\fill}
   \end{center}
        
 % Isso na versao final do trabalho!!!       
   Trabalho aprovado. \imprimirlocal, 01 de janeiro de 2014:

   \assinatura{\textbf{\imprimirorientador} \\ Orientador} 
   \assinatura{\textbf{\imprimircoorientador} \\ Co-Orientador} 
   \assinatura{\textbf{Professor} \\ Convidado 1}
   \assinatura{\textbf{Professor} \\ Convidado 2}
   \assinatura{\textbf{Professor} \\ Convidado 3}
      
   \begin{center}
    \vspace*{0.5cm}
    {\large\imprimirlocal}
    \par
    {\large\imprimirdata}
    \vspace*{1cm}
  \end{center}
  
\end{folhadeaprovacao}
% ---

% Dedicatória
% ---
% Dedicatória
% ---
\begin{dedicatoria}
   \vspace*{\fill}
   \centering
   \noindent
   \textit{ Aos verme que roeu as frias carnes de meu cadáver.} \vspace*{\fill}
\end{dedicatoria}
% ---

% Agradecimentos
% ---
% Agradecimentos
% ---
\begin{agradecimentos}



Agradeço a Xuxa, meus pais, cachorro, gato e papagaio, por ...

Agradeço ao meu orientador, XXXXXXXXX, por todos os conselhos, pela paciência e ajuda nesse período.

Aos meus amigos ...

Aos professores ...

À XXXXXX pelo apoio financeiro para realização deste trabalho de pesquisa.

\end{agradecimentos}
%% ---

% Epígrafe
% ---
% Epígrafe
% ---
\begin{epigrafe}
    \vspace*{\fill}
	\begin{flushright}
		\textit{``Não sei o que, \\
		          não sei o que,\\
                  não sei o que lá.''\\
		          (Autor Desconhecido)}
	\end{flushright}
\end{epigrafe}
% ---

% Resumo e Abstract
% ---
% RESUMOS
% ---

% RESUMO em português
\setlength{\absparsep}{18pt} % ajusta o espaçamento dos parágrafos do resumo
\begin{resumo}
 Segundo a ABNT, o resumo deve ressaltar o
 objetivo, o método, os resultados e as conclusões do documento. A ordem e a extensão
 destes itens dependem do tipo de resumo (informativo ou indicativo) e do
 tratamento que cada item recebe no documento original. O resumo deve ser
 precedido da referência do documento, com exceção do resumo inserido no
 próprio documento. Umas 10 linhas (\ldots) As palavras-chave devem figurar logo abaixo do
 resumo, antecedidas da expressão Palavras-chave:, separadas entre si por
 ponto e finalizadas também por ponto.

 \textbf{Palavras-chaves}: lixo eletrônico. jogo. jogo educativo. computação. ecologia. resíduo eletrônico. REEE. e-lixo.
\end{resumo}

% ABSTRACT in english
\begin{resumo}[Abstract]
 \begin{otherlanguage*}{english}
   This is the english abstract.

   \vspace{\onelineskip}
 
   \noindent 
   \textbf{Keywords}: latex. abntex. text editoration.
 \end{otherlanguage*}
\end{resumo}

% Lista de ilustrações
\pdfbookmark[0]{\listfigurename}{lof}
\listoffigures*
\cleardoublepage

% Lista de tabelas
\pdfbookmark[0]{\listtablename}{lot}
\listoftables*
\cleardoublepage

% Lista de abreviaturas e siglas
\begin{siglas}
  \item[ABNT] Associação Brasileira de Normas Técnicas
  \item[abnTeX] Normas para TeX
\end{siglas}

% Lista de símbolos
\begin{simbolos}
  \item[$ \Gamma $] Letra grega Gama
  \item[$ \Lambda $] Lambda
  \item[$ \zeta $] Letra grega minúscula zeta
  \item[$ \in $] Pertence
\end{simbolos}

% Inserir o SUMÁRIO
\pdfbookmark[0]{\contentsname}{toc}
\tableofcontents*
\cleardoublepage

% ----------------------------------------------------------
% ELEMENTOS TEXTUAIS (Capítulos)
% ----------------------------------------------------------
\textual
% Elementos textuais com numeração arábica
\pagenumbering{arabic}
% Reinicia a contagem do número de páginas
\setcounter{page}{1}

% Inclui cada capitulo da Dissertação
% ----------------------------------------------------------
% Introdução 
% Capítulo sem numeração, mas presente no Sumário
% ----------------------------------------------------------

\chapter*[Justificativa]{Justificativa}
\addcontentsline{toc}{chapter}{Justificativa}

Além do desperdício de resíduos gerados pela sociedade humana é um problema a ser enfrentado, desde a quarta revolução[1] e o massivo número de produtos eletrônicos sendo produzidos e descartados, passou a ser um problema o que fazer com esses resíduos novos, pilhas, baterias, celulares, televisores e computadores por exemplo. Existe a preocupação maior com estes descartados, pois em sua vez, possuem componentes internos como baterias e outros elementos que podem ser extremamente danosos ao meio ambiente, além de possuir metais preciosos e outros elementos raros e que poderiam ser reaproveitados na manufatura de outros componentes.

Mas como mudar esse problema? Podemos incentivar as pessoas por meio da sensibilização e conscientização por meio da educação[2] mostrou um ganho na compreensão e mudanças de hábitos. Existe uma proposta de jogo eletrônico baseado em Novas Tecnologias da Informação e Comunicação, como uma ferramenta para auxiliar professores para abordar este tipo de assunto de uma forma mais interativa com os seus alunos.[3] Seguindo este exemplo: What Weee Are, é um jogo eletrônico que aborda o assunto dos resíduos eletrônicos e pode ser usado como uma ferramenta na conscientização das pessoas quanto ao assunto de resíduos eletrônicos.


\section*{Motivação}\label{sec:motivacao}
\addcontentsline{toc}{section}{Motivação}

O motivo de utilizar um jogo como motor para impulsionar o interesse dos mais jovens para questões ambientais como essa, parte de que os jogos estão amplamente presentes na vida das pessoas, e podem desenvolver habilidades de cooperação, colaboração em desafios para resolver problemas[4]. Podendo ser uma ferramenta utilizada pelos professores para fomentar o desenvolvimento cognitivo dos alunos, pois mostra uma ferramenta capaz de motivar e despertar o interesse à diversos assuntos.[5]

\section*{Objetivos}\label{sec:objetivos}
\addcontentsline{toc}{section}{Objetivos}

\lipsum[36]


% PARTE - Define a divisão do documento em partes (Não é obrigatório)
\part{WhatWeeeAre: jogo eletrônico com abordagem à reciclagem de resíduos eletrônicos}
% ---
\chapter{WhatWeeeAre: jogo eletrônico com abordagem à reciclagem de resíduos eletrônicos}
% ---

What Weee Are é um projeto multimídia, um jogo eletrônico que aborda os desperdício de recursos de nossa sociedade, a ideia é que o jogo seja capaz de ser usado como uma ferramenta educacional utilizada pelos professores, capaz de motivar o interesse dos alunos a um dos problemas atuais de nossa sociedade. Utilizar um jogo como ferramenta de educação já foi apresentado como viável e devemos aproveitar as oportunidades que o atual mundo tecnológico nos apresenta[5]. 

What Weee Are é um jogo de plataforma 2D, como elementos de RPG para progressão do personagem, sendo capaz de coletar recursos chaves(ouro, cobre, ferro e plantas) que podem ser utilizados para comprar aprimoramentos ao personagem. Outra maneira de conseguir recursos é através da coleta de materiais não brutos como: cabos, placas, mouse. E desmontá-los a fim de obter seus recursos brutos. Utilizando fotografias reais como material gráfico para compor as cenas dos jogos, personagem, recursos, obstáculos.

%-
\section{Metodologias}
%-

Para o desenvolvimento do projeto, foi utilizado o motor gráfico UNITY, disponibilizado gratuitamente em https://unity.com/pt, com auxilio da linguagem de programação C\# para criação de scripts e utilitários ao jogo.

\subsection{Características necessárias}
Durante a concepção do jogo, foram analisados os conceitos de jogos de plataformas que fizeram sucesso e tomadas como inspiração para o desenvolvimento, como a franquia Donkey Kong Country e Super Mario Bros, ambas pertencentes à Nintendo e jogos de RPG onde existem uma progressão do personagem e novas habilidades adquiridas, como a série The Legend of Zelda, Final Fantasy:

\begin{itemize}
    \item Tipos diferentes de pulo, como por exemplo segurar o botão de pular e efetuar um salto longo.
    \item Coleta de itens e sistema de Crafting
    \item Habilidades desbloqueáveis
    \item Plataformas dinâmicas
    \item Inimigos de diferentes aspectos
    \item Movimentação realizada pelo jogador
\end{itemize}

\subsection{Objetivos do jogo}
O Objetivo de What Weee Are é o jogador passar pelas 4 fases do jogo. Em cada fase o jogador será confrontado com um número de inimigos e itens a serem coletados. Nas fases estarão a disposição as salas de Assembler e Disassembler para auxiliá-lo na obtenção de itens complexos e novas habilidades. Com auxílio do Grilo Falante, o mentor do personagem, será apresentada a narrativa do jogo e explicando os eventos do universo a cada nova fase encontrada. 

Em cada fase, os desafios serão incrementados e necessitará que o jogador consiga as habilidades especiais disponíveis pelo sistema de Crafting, para acessar pontos em que sem tais habilidades se tornará impossível, forçando com que o jogador interaja com o ambiente e colete o maior número de itens possíveis nas fases. 

Há dois tipos de inimigos nas fases, as Formigas, e as Aranhas, e suas variantes, que possuem características diferentes e dificuldades incrementadas a cada nova fase. Quando o jogador morre (atinge o limite mínimo de vidas) ele retorna ao início da fase.


\subsection{Desenvolvimento}
Durante todo o processo de desenvolvimento, foram efetuados testes e questionamentos sobre conceitos que são necessários para incrementar a experiência do jogador.

Inicialmente o projeto foi iniciado desenvolvendo mecânicas que poderiam ser reutilizadas ao longo de todo o progresso do jogo, mecânicas essas que foram sendo aprimoradas durante o desenvolvimento e conforme era necessário alguma nova característica se tornasse necessária para progredir. 

Inicialmente foram desenvolvidas as interações do personagem Weee com o ambiente e os inimigos, interações essas de morte ao chegar à 0(zero) vidas disponíveis, dano ao chocar-se com inimigos pelas laterais e por cima, e podendo derrotar os inimigos ao pular sobre eles. Nessas interações com os inimigos, ao tomar dano o inimigo será levemente empurrado na direção em que fora atingido pelo Weee, o mesmo se aplica ao Weee ao ser atingido pelos inimigos, essa mecânica se chama knockback.

Assim fora iniciado o desenvolvimento do sistema de coleta de itens e de inventário do personagem, onde seria possível visualizar a quantidade de diferentes itens coletados. 

<imagem 1>

Os itens coletados podem ser refinados em itens mais complexos ou desmembrados em itens mais primitivos, esta definição pode ser efetuada a partir de como o item é produzido, este sistema de receitas em conjunto aos sistemas de montagem e desmontagem, que irão consumir os itens produtores para a criação dos itens resultantes:

<imagem 2>\\
<imagem 3>\\
<imagem 4>\\

Para tudo isso poder funcionar e ter a dinâmica de levar os itens aos compartimentos para montá-los e desmontá-los, foi utilizado o conceito de DragAndDrop onde o usuário poderá clicar e arrastar os itens até os compartimentos de entrada.

% ---
\subsection{Tecnologias utilizadas}
Foi utilizado a engine para jogos Unity, versão 2021, focado nas ferramentas em 2D, como:
\begin{itemize}
    \item Tile Pallet: ferramenta de paleta para aplicar os Tiles (imagens/sprites) no jogo 2D, dentro do cenário, utilizando um pincel é permitido pintar obstáculos e plataformas de maneira facilitada ao desenvolvimento.
    \item Tile rule: ferramenta auxiliar ao tile pallet, que ajuda a criar uma regra de quais sprites devem ficar seguindo um ao outro e montar um cenário coerente com bordas, segmentos, rampas, etc.
    \item Animators: Ferramenta auxiliar para criar animações dentro do jogo, como animação de andar, pular, morrer.
    \item Localization: Ferramenta disponibilizada pela Unity, que auxilia na localização do jogo em diversos idioma, possibilitando assim o desenvolvedor a promover a localização do jogo em diversas linguas com mais facilidades, por meio de tabelas de texto e também tabelas de objetos, podendo conter texturas por exemplo.
\end{itemize}

\subsection{Implementação}
\lipsum[34-48]





\chapter{\textit{Storyline}}\label{cap:estArte}

\lipsum[34]

\section*{Nível 1: A cozinha}\label{sec:primTrab}
\addcontentsline{toc}{section}{Nível 1: A cozinha}

\lipsum[34-36]

\section*{Nível 2: O Jardim}\label{sec:primTrab}
\addcontentsline{toc}{section}{Nível 2: O Jardim}

\lipsum[34-36]

\section*{Nível 3: O Metrô}\label{sec:primTrab}
\addcontentsline{toc}{section}{Nível 3: O Metrô}

\lipsum[34-36]

\section*{Nível 4: O Vulcão}\label{sec:primTrab}
\addcontentsline{toc}{section}{Nível 4: O Vulcão}

\lipsum[34-36]
\chapter{Componentes}\label{cap:ferramentas}

\lipsum[10]

\section{Inímigos}
\lipsum[23]
\begin{itemize}
    \item inimigo1
    \item inimigo2
\end{itemize}

\section{Itens}
\lipsum[23]
\begin{itemize}
    \item item1
    \item item2
\end{itemize}

\section{Melhorias}
\lipsum[23]
\begin{itemize}
    \item Pulo duplo
    \item Velocidade aumentada
    \begin{itemize}
        \item 3x ferro
        \item 2x ouro
    \end{itemize}
\end{itemize}

\section{Aliádos}
O Grilo falante

\lipsum[23]

\chapter{Sistema Proposto}\label{cap:proposta}

Esse trabalho propõe um sistema de... 


\section{Primeira Parte do Sistema Proposto}

\lipsum[67]

\section{Considerações Finais}

\lipsum[68]

\chapter{Resultados e Discussão}\label{cap:resultados}

\lipsum[73]

\section{Base de Dados}

\lipsum[72]

\section{Considerações Finais}

\lipsum[74]
\chapter*{Conclusões e Trabalhos Futuros}\label{cap:conclusao}
\addcontentsline{toc}{chapter}{Conclusão e Trabalhos Futuros}

\lipsum[81]

\section*{Conclusões}

\lipsum[82-84]

\section*{Trabalhos Futuros}

\lipsum[85] 

% ----------------------------------------------------------
% ELEMENTOS PÓS-TEXTUAIS (Referências, Glossário, Apêndices)
% ----------------------------------------------------------
\postextual

% Referências bibliográficas
\bibliography{bibliografia}

% Glossário (Consulte o manual)
%\glossary

% Apêndices
% ----------------------------------------------------------
% Apêndices
% ----------------------------------------------------------

% ---
% Inicia os apêndices
% ---
\begin{apendicesenv}

% Imprime uma página indicando o início dos apêndices
\partapendices

% ----------------------------------------------------------
\chapter{Primeiro Apêncice}
% ----------------------------------------------------------

\lipsum[50] % Texto qualquer. REMOVER!!

% ----------------------------------------------------------
\chapter{Segundo apêndice com título tão grande quanto se queira porque ele já faz a quebra de linha da coisa toda}
% ----------------------------------------------------------
\lipsum[51-53] % Texto qualquer. REMOVER!!

\end{apendicesenv}
% ---

% Anexos
% ----------------------------------------------------------
% Apêndices
% ----------------------------------------------------------

% ---
% Inicia os anexos
% ---
\begin{anexosenv}

% Imprime uma página indicando o início dos anexos
\partanexos

% ---
\chapter{Nome do Primeiro Anexo}
% ---
\lipsum[30] % Texto qualquer. REMOVER!!

% ---
\chapter{Nome de Outro Anexo}
% ---

\lipsum[32] % Texto qualquer. REMOVER!!

\end{anexosenv}

% Índice remissivo (Consultar manual)
%\phantompart
%\printindex

\end{document}
